% ======================================================================
% Проверка реализации (таблица). Ведёт D — Филиппов Андрей.
% Каждый владелец присылает строки за свой блок (PR меняет только свою строку).
%
% Требования (main.tex, строки 1885–1891): столбцы «Задача», «Команда запуска»,
% «Проверяемая величина», «Результат проверки»; включить сравнение интегралов
% и объёмов с аналитическими значениями и проверки численных методов из п. 3.5.
% В результате — конкретное расхождение или выполненное условие.
% Полный журнал вывода программы вместо таблицы не приводится.
% Образец: report_template.tex, строки 313–336.
% ======================================================================

\begin{table}[htbp]
\centering\small
\caption{Проверка реализации}
\label{tab:checks}
\begin{tabular}{|p{1.4cm}|p{3.9cm}|p{4.6cm}|p{4.6cm}|}\hline
Задача & Команда запуска & Проверяемая величина & Результат проверки\\\hline
1.1 & \texttt{python experiments.py 1.1} & Интеграл $I_1$: расхождение с аналитическим значением на сетках $4,\dots,64$ (оба метода) & —\\\hline
1.2 & \texttt{python experiments.py 1.2} & Интеграл $I_2$: расхождение с аналитическим значением на сетках $4,\dots,64$ (оба метода) & —\\\hline
2.1 & \texttt{python experiments.py 2.1} & Интеграл $I_{\mathrm{п}}$ и объём: расхождение с аналитикой и $V=abc$ & —\\\hline
2.2 & \texttt{python experiments.py 2.2} & Интеграл $I_{\mathrm{ц}}$ и объём: расхождение с аналитикой и $V=\pi R^2H/2$ & —\\\hline
2.3 & \texttt{python experiments.py 2.3} & Интеграл $I_{\mathrm{с}}$ и объём: расхождение с аналитикой и $V=2\pi R^3/3$ & —\\\hline
3.1–3.4 & \texttt{python experiments.py 3.1} и т.д. & Значения формул: расхождение аналитического и программного расчёта & —\\\hline
3.5.1 & \texttt{python experiments.py 3.5.1} & Оптимальные шаги разностей; минимум ошибки на найденном шаге & —\\\hline
3.5.2 & \texttt{python experiments.py 3.5.2} & Ошибки прогнозов $P_1$–$P_3$ (п. 3.2.3) от $h$: рост при $h\to s$ & —\\\hline
3.5.3 & \texttt{python experiments.py 3.5.3} & Граница устойчивости по $\Delta t$: сравнение с линейным критерием & —\\\hline
3.5.4 & \texttt{python experiments.py 3.5.4} & Пороги флаппинга: сравнение с $k_*$; устойчивость наблюдения при уменьшении шага & —\\\hline
3.5.5 & \texttt{python experiments.py 3.5.5} & Координаты минимума ($\Delta_c$); допустимость точек; невозрастание затрат & —\\\hline
3.6 & \texttt{python experiments.py 3.6} & Средние затраты: сходимость при уменьшении $\Delta t$ & —\\\hline
\end{tabular}
\end{table}
